\documentclass[11pt,a4paper]{article}

\usepackage[T1]{fontenc}
\usepackage[utf8]{inputenc}
\usepackage{newtxtext,newtxmath}
\usepackage[margin=2.2cm]{geometry}
\usepackage{graphicx}
\usepackage{booktabs}
\usepackage{amsmath}
\usepackage[hidelinks]{hyperref}
\usepackage{url}
\usepackage[numbers,sort&compress]{natbib}
\usepackage{caption}
\captionsetup{font=small,labelfont=bf}
\usepackage{enumitem}
\setlist{nosep,leftmargin=*}

\graphicspath{{figures/}}

\title{\textbf{Experimental Sensitivity Analysis of Kubernetes Horizontal Pod Autoscaler Configurations under Steady, Bursty, and Periodic Web Workloads}}

\author{
  Basim Kafeel\\
  \small 22I-0775 --- \texttt{i220775@nu.edu.pk}\\
  \small Primary Researcher
  \and
  Ali Aamir Ashraf\\
  \small 22I-1199 --- \texttt{i221199@nu.edu.pk}\\
  \small Co-Researcher
  \\[1ex]
  \small Department of Computer Science, FAST-NUCES Islamabad\\
  \small BS Computer Science, Section CS-B, Group G01
}
\date{}

\begin{document}
\maketitle

%% ABSTRACT (written after the results are final)

\section{Introduction}

Web applications rarely receive a constant amount of traffic. Demand can stay
roughly level, jump sharply for a short time, or rise and fall in a regular
cycle. Provisioning for the peak wastes capacity for most of the day, while
provisioning for the average lets response times grow when demand rises.
Elasticity, the ability to add and remove resources as demand changes, is
therefore one of the main reasons for running web services in the
cloud~\cite{aldhuraibi2018elasticity,qu2019autoscaling}. Deciding \emph{when}
and \emph{how much} to scale is the job of an autoscaler, and decades of work
have produced rule-based, queueing-based, control-theoretic, time-series and
learning-based
approaches~\cite{loridobotran2014a,qu2019autoscaling,chen2019a,gari2021reinforcement}.

Containers and container orchestration have made autoscaling a standard
platform feature rather than custom
infrastructure~\cite{bernstein2014containers,pahl2019cloud,burns2016borg}.
Kubernetes, now the dominant orchestrator~\cite{truyen2019a,carrion2023kubernetes},
ships a Horizontal Pod Autoscaler (HPA) that periodically compares an observed
metric, most often CPU utilisation, with a target and changes the number of
pod replicas accordingly~\cite{nguyen2020horizontal}. The HPA is a closed
feedback loop that monitors the application, decides on a replica count and
acts on it, which places it in the area of autonomic and self-adaptive
systems~\cite{kephart2003the,salehie2009selfadaptive}.

Although many studies propose new autoscalers that replace or extend the
HPA~\cite{dangquang2021deep,imdoukh2020machine,toka2021machine,phuc2022trafficaware,abdelkhaleq2023intelligent},
the standard HPA remains the most widely used option because it is built in
and easy to operate~\cite{pozdniakova2024slaadaptive}. Its behaviour, however,
depends on several parameters: the utilisation target, the replica bounds, and
since the \texttt{autoscaling/v2} API, scale-up and scale-down policies and
stabilisation windows. These parameters are usually left at their defaults or
set by trial and error. Evaluations of autoscalers have also been criticised
for testing a single scenario and comparing only against static
provisioning~\cite{ilyushkin2018an,papadopoulos2016peas}.

This article asks a practical question: \emph{how do different Kubernetes HPA
configurations affect application performance, SLA compliance and resource
efficiency under steady, bursty and periodic web workloads?} We answer it with
a controlled, repeated experiment on a local Kubernetes cluster. The same
CPU-bound web service is run under four configurations (fixed provisioning,
the default HPA, a responsive HPA and a balanced HPA) and three workload
patterns, with three repetitions of each combination. The contributions are:
\begin{itemize}
  \item a fully scripted and reproducible experimental setup (application,
        manifests, load generator, runner and analysis) published as open
        source;
  \item a measurement of how the HPA's target and \texttt{behavior} settings
        trade latency, SLA violations and scaling activity against
        provisioned capacity (pod-seconds) for three workload shapes; and
  \item practical recommendations, grounded in the measured data, for choosing
        an HPA configuration for a given traffic pattern.
\end{itemize}

The rest of the article is organised as follows. Section~\ref{sec:background}
explains how the HPA works. Section~\ref{sec:related} reviews related work.
Section~\ref{sec:method} describes the experimental design.
Section~\ref{sec:results} presents the results, Section~\ref{sec:discussion}
discusses them together with threats to validity, and
Section~\ref{sec:conclusion} concludes.

\section{Background: the Kubernetes Horizontal Pod Autoscaler}
\label{sec:background}

A Kubernetes \emph{Deployment} keeps a requested number of identical pods
running, and a \emph{Service} spreads incoming connections over the pods that
are ready. The HPA controller adjusts the Deployment's replica count. Every
synchronisation period (15\,s by default) it reads the current metric for the
target pods from the Metrics API, which the Metrics Server fills by scraping
the kubelets, and computes
\begin{equation}
  \text{desiredReplicas} = \left\lceil \text{currentReplicas} \times
  \frac{\text{currentMetricValue}}{\text{targetMetricValue}} \right\rceil .
  \label{eq:hpa}
\end{equation}
For CPU, the metric is average utilisation expressed as a percentage of each
pod's CPU \emph{request}, so a pod that requests 250 millicores and uses
200 millicores is at 80\%. Changes are skipped when the ratio is within a
tolerance of 10\%, and the result is clamped to the configured
\texttt{minReplicas} and \texttt{maxReplicas}~\cite{nguyen2020horizontal}.

The \texttt{autoscaling/v2} API adds a \texttt{behavior} section with separate
rules for scaling up and down. Each direction has \emph{policies} that cap the
change per period (for example, at most four pods or 100\% of the current
count per 15\,s), a \texttt{selectPolicy} that picks the most or least
permissive policy, and a \emph{stabilisation window}. During a window the
controller uses the most conservative recommendation it has computed, the
highest for scale-down and the lowest for scale-up, which damps
oscillation. When no behaviour is given, the defaults are: scale-up with no
window and up to the larger of 100\% or four pods per 15\,s; scale-down with a
300\,s window and up to 100\% per 15\,s.\footnote{Kubernetes documentation,
``Horizontal Pod Autoscaling'',
\url{https://kubernetes.io/docs/concepts/workloads/autoscaling/horizontal-pod-autoscale/}.}
Because new pods need time to start, pass readiness checks and appear in the
Metrics API, the effective reaction time of the loop is longer than the
synchronisation period. Recent work analyses the stability of this control
loop formally~\cite{serracanta2025on}.

\section{Related Work}
\label{sec:related}

\paragraph{Autoscaling surveys.}
Lorido-Botran et al.\ classify cloud autoscaling techniques into threshold
rules, reinforcement learning, queueing theory, control theory and time-series
analysis~\cite{loridobotran2014a}. Qu et al.\ survey autoscaling specifically
for web applications and organise the problem along monitoring, analysis,
planning and execution~\cite{qu2019autoscaling}. Al-Dhuraibi et al.\ review
elasticity in cloud computing, including container
elasticity~\cite{aldhuraibi2018elasticity}, and Chen et al.\ frame autoscalers
as self-aware and self-adaptive systems~\cite{chen2019a}. More recent surveys
cover autoscaling of containers in cloud and edge/fog
settings~\cite{dogani2023autoscaling}, reinforcement-learning
autoscalers~\cite{gari2021reinforcement} and machine-learning-based container
orchestration~\cite{zhong2022machine}. Container technologies and orchestrators
themselves are reviewed in~\cite{pahl2019cloud,truyen2019a}, and Kubernetes
scheduling in~\cite{carrion2023kubernetes,senjab2023a}.

\paragraph{Classic dynamic provisioning.}
Before containers, dynamic provisioning of multi-tier Internet applications
combined queueing models with predictive and reactive
control~\cite{urgaonkar2008agile}, and AutoScale showed that capacity can be
added and removed robustly under unpredictable load while meeting
response-time targets~\cite{gandhi2012autoscale}. Brownout, which degrades
optional content under overload, is a complementary
approach~\cite{xu2020brownout}. Workload prediction with ARIMA models was used
to provision ahead of demand~\cite{calheiros2015workload}, and Chameleon
combines several forecasting methods in a hybrid proactive
autoscaler~\cite{bauer2019chameleon}.

\paragraph{Studies of the Kubernetes HPA.}
Nguyen et al.\ study HPA behaviour experimentally and compare Kubernetes
resource metrics with Prometheus custom metrics, offering lessons for tuning
the HPA~\cite{nguyen2020horizontal}. Casalicchio compares relative and absolute
CPU measures for autoscaling CPU-intensive containerised
applications~\cite{casalicchio2019a}. Taherizadeh and Stankovski propose
multi-level rules that combine infrastructure and application
metrics~\cite{taherizadeh2019dynamic}. Pozdniakova et al.\ adjust the HPA's
utilisation threshold from data so that a response-time SLO is met with little
over-provisioning~\cite{pozdniakova2024slaadaptive}, and Augustyn et al.\ tune
the maximum pod count of individual HPAs from a load profile, reducing the
number of nodes in a medical-records service~\cite{augustyn2024tuning}.
Serracanta et al.\ examine the stability of the HPA control
loop~\cite{serracanta2025on}. For edge clusters, traffic-aware and node-based
HPA variants distribute pods according to where requests
arrive~\cite{phuc2022trafficaware,phuc2022nodebased}.

\paragraph{Proactive and learning-based alternatives.}
Several works replace the reactive HPA rule with a forecast or a learned
policy: a bidirectional LSTM~\cite{dangquang2021deep}, neural-network
forecasting~\cite{imdoukh2020machine}, machine-learning-based scaling for edge
clusters~\cite{toka2021machine,toka2020adaptive}, reinforcement
learning~\cite{abdelkhaleq2023intelligent} and adaptive microservice
scaling~\cite{cruzcoulson2020adaptive}. LARE-HPA co-optimises latency and
resource efficiency for horizontal pod autoscaling~\cite{kim2025larehpa}.
Container placement and cost-efficient orchestration in Kubernetes are studied
in~\cite{rossi2020geodistributed,zhong2020a}.

\paragraph{Evaluation methodology.}
Ilyushkin et al.\ point out that new autoscalers are seldom compared with the
state of the art and evaluate seven policies across operating
conditions~\cite{ilyushkin2018an}. PEAS evaluates autoscalers under workload
uncertainty and shows that conclusions from single scenarios do not
generalise~\cite{papadopoulos2016peas}. Papadopoulos et al.\ give
methodological principles for reproducible performance evaluation in the
cloud, such as repeated runs and reporting of
variability~\cite{papadopoulos2021methodological}, and Leitner and Cito show
how much performance varies in public clouds~\cite{leitner2016patterns}.
Formal performance models have also been used to verify autoscaling
policies~\cite{evangelidis2018performance}, and uncertainty in elasticity
controllers is discussed in~\cite{jamshidi2016managing}. Tail latency, rather
than the mean, often determines user experience at
scale~\cite{dean2013the}, which is why we report the 95th percentile.

\paragraph{Gap.}
Most of the work above proposes a new autoscaler or tunes a single parameter
for a specific service. Few studies compare the \emph{built-in} HPA settings,
the CPU target together with the \texttt{behavior} policies, against each
other and against fixed provisioning across several workload shapes, with
repeated runs and openly available code and raw data. This study fills that
gap for a controlled, CPU-bound web service.

%% METHOD, RESULTS, DISCUSSION, CONCLUSION follow once the data are final.

\bibliographystyle{unsrtnat}
\bibliography{references}

\end{document}
